\subsection{Physical Edges of the Domains}
\label{sec:domain-boundaries}

Each map ends where its stable solutions end, and the end is a curve, not a rectangle. Scalarized black holes exist only above a smallest horizon radius $r_{h,\min}(\beta)$, so the $(\beta, \rh)$ domain is bounded from below: over the sampled couplings $\beta$ from $\bhEdaEdgeBetaLow$ to $\bhEdaEdgeBetaHigh$ the edge runs from $\bhEdaEdgeRhMinLow$ to $\bhEdaEdgeRhMinHigh$ (Section~\ref{sec:eda-bh}). A neutron-star sequence of fixed $(\beta, \lambda)$ is stable only up to its maximum mass $M_{\max}$, the turning point along $\rhoc$; beyond it the branch is unstable (Section~\ref{sec:eda-ns}).

These edges, with the sampled coupling ranges, are where the map ends physically, and so where a prediction has no meaning. They are the domain the surrogate is guarded by (Section~\ref{sec:validity}).
